\subsection{The complete interactive prover and verifier}\label{sec:full:protocol}

\begin{definition}[Explicit relation-proof interface]\label{def:full:proof-engine}
Fix an interactive zero-knowledge \emph{proof-of-relation}
engine $\Pi_{\rm rel}$ for polynomial-size arithmetic and
Boolean circuits over $\F$, including:
valid field-symbol encodings, Pedersen bit openings,
canonical Merkle-tree reconstruction, polynomial
arithmetic and equality, and all public constraints
specified below.
Its public input is a tuple of public commitments,
roots, points and claimed scalars. Its private
witness is the collection of field values,
blinding scalars, word evaluations and paths
required to evaluate the circuit.
The verifier checks the \emph{entire} relation
without receiving any private circuit input.

We separate two possible engines:
\begin{enumerate}[label=\textup{(\roman*)}]
\item $\Pi_{\rm full}$ certifies the entire
committed word at each level, its polynomial
degree bound, source relation and fold
identity (a nonsuccinct reference mode).
\item $\Pi_{\rm sparse}$ certifies just
the explicitly requested authenticated
local statements; its soundness requires
a separately proved proximity theorem.
\end{enumerate}
The existence of a complete relation
circuit is an algorithmic statement,
not a claim that the source KBFold
proximity theorem transfers to the
sparse mode.
\end{definition}

\begin{definition}[Public commitment registry]\label{def:full:registry}
For each named hidden field variable, keep
a single public vector of Pedersen bit
commitments. An invocation of
$\Pi_{\rm rel}$ must refer to that same
vector whenever it uses the variable.
At preprocessing the registered variables
include all entries $f(a)$, all
coefficients of $R$, and every
symbol of the level-zero masked word.
At each Sumcheck round the registered
variables include the three coefficients
of $\mu_j$, the round mask $c_j$
and, after $r_{j-1}$, the offset $b_j$
and every newly committed level-$j$
oracle symbol.
The equality of a named field variable
across two circuit invocations is checked
through its fixed public commitment vector.
\end{definition}

\begin{definition}[Complete preprocessing relation]\label{def:full:pre-relation}
The stage-zero public statement consists
of the initial Merkle root
$\mathsf{root}_0$, the public source
and mask commitment vectors
$\mathbf C_f,\mathbf C_R$, and parameters.
The private witness of
$\mathcal R_{\rm pre}$ consists of
$(f,R)$, every related blinding vector
and the full level-zero Merkle tree.
The relation checks:
\begin{enumerate}[label=\textup{(\roman*)}]
\item every source/mask commitment
opens to the correctly encoded
field value, with $\deg R<B$;
\item $Q_0=U_f^{(m)}+X^NR$;
\item every level-zero leaf
commitment opens to $Q_0(x)$
and all its hashes reconstruct
the announced $\mathsf{root}_0$;
\item $\deg Q_0<D_0$.
\end{enumerate}
The preprocessing proof is a
\emph{knowledge-binding registration
proof} completed before $z$ is
chosen. Its extraction property
is a named cryptographic interface
contract of Chapter~14.
\end{definition}

\begin{definition}[Round relations]\label{def:full:round-relations}
For round $1\le j\le m$, the
\emph{pre-challenge relation}
$\mathcal R_{j,\rm pre}$, proved
after commitment of $\mu_j$
but before sampling $r_{j-1}$,
binds the hidden $c_j$,
the three coefficients of
$\mu_j$, and earlier offset
$b_{j-1}$ to
\[
 \mu_j(T)=s_j(T)+
   b_{j-1}/2+c_j(2T-1),
\]
where $s_j$ is reconstructed
from the same committed source
table $f$ and already fixed
earlier challenges.

The \emph{post-challenge relation}
$\mathcal R_{j,\rm post}$, proved
after $r_{j-1}$ but before any
level-$j$ query, checks
\[
 \begin{aligned}
 b_j&=b_{j-1}/2+c_j(2r_{j-1}-1),\\
 \mu_j(0)+\mu_j(1)&=\sigma_{j-1}+b_{j-1},\\
 \mu_j(r_{j-1})&=\sigma_j+b_j,\\
 Q_j&=\pfold_{r_{j-1}}(Q_{j-1}),\\
 \deg Q_j&<D_j,
 \end{aligned}
\]
and, in $\Pi_{\rm full}$ mode,
checks all openings of the
level-$j$ Merkle leaf commitments
against the coefficients of $Q_j$.
The public challenge is not
allowed to modify any earlier
commitment or named hidden value.
\end{definition}

\begin{definition}[Full protocol $\Pi_{\rm KBFold}^{\rm hid}$]\label{def:full:protocol}
The interactive opening protocol is:

\begin{enumerate}[label=\textup{(\roman*)}]
\item \textbf{Before evaluation:}
the prover samples $R$ and fresh
bit-commitment blindings, builds
$Q_0$, and publishes
$(\mathbf C_f,\mathbf C_R,\mathsf{root}_0)$.
It proves $\mathcal R_{\rm pre}$.
The verifier verifies this proof.
No evaluation point $z$ has yet
been chosen.
\item \textbf{Evaluation claim:}
after preprocessing, the verifier
chooses $z\in\F^m$, and the
prover publishes a claimed scalar
$v$. An honest prover sends
$v=\widetilde f(z)$.
The prover initializes $\chi_0$,
the claimed running sum
$\sigma_0=v$, and $b_0=0$.
The relation
$v=\sum_a f(a)\chi_0(a)$
against the \emph{preprocessing-bound}
source table is included in the
first round's proof.
A cheating prover may choose
a false $v$; the verifier
must still enforce the relation.
\item \textbf{For $j=1,\ldots,m$:}
the prover computes $s_j$ and
samples fresh $c_j$, commits to
the three coefficients of $\mu_j$
and $c_j$, and proves
$\mathcal R_{j,\rm pre}$.
Only then the honest verifier
samples the fresh fold/Sumcheck
challenge $r_{j-1}\in\F$.
The prover computes $b_j,Q_j$,
publishes their fresh hiding
commitments and $\mathsf{root}_j$,
and proves $\mathcal R_{j,\rm post}$.
\item \textbf{Authenticated queries:}
after $\mathsf{root}_j$ is bound,
the verifier may choose up to
$q_j$ adaptive indices
$y\in\Omega_j$ for the previous
fold $j-1$. For each it receives
the three shells of
Definition~\ref{def:full:fold-query}
and a zero-knowledge proof of
their hidden fold equation and
registered leaf openings.
It checks every Merkle path
and proof. No field-symbol
value is publicly revealed.
\item \textbf{Final closure:}
after all $m$ rounds, the
verifier requests two distinct
final-level authenticated shells
and a proof of the
degree-$<2$ final relation and
the terminal Sumcheck link of
Definition~\ref{def:full:final-close}.
It accepts if and only if
all public Merkle checks and
all required relation proofs
have accepted.
\end{enumerate}
The full verifier view records
every commitment vector and root,
all inner proof messages,
$z,v,r$, adaptive query indices,
all authenticated shells and
all final proof messages.
\end{definition}

\begin{lemma}[All protocol relations are finite arithmetic circuits]\label{lem:full:circuits}
For fixed finite field representation,
$m,B,M$ and a bounded number of queries,
each relation in Definitions
\ref{def:full:pre-relation},
\ref{def:full:round-relations},
\ref{def:full:fold-query} and
\ref{def:full:final-close}
is decidable by an explicit finite
Boolean/arithmetic circuit of
size polynomial in the expanded
input and committed-word lengths
and the hash-circuit complexity.
\end{lemma}
\begin{proof}
A field element has a finite
canonical $b$-bit encoding.
Every field sum, product, inverse
of a fixed nonzero public scalar
and basis expansion has a finite
arithmetic circuit over the fixed
representation. The polynomials
$U_f,Q_j,s_j,H_j,\mu_j$ have
publicly bounded degrees and
are computed by finitely many
sums and products with indices
bounded by $N,B,M$. A Merkle
leaf hash and its internal
hash circuits are finite by
assumption on the efficiently
computable hash $H$; the
bounded Merkle reconstruction
is a composition of such
circuits. Group scalar
exponentiation to check a
Pedersen commitment is
computable by repeated
squaring in $O(\log P)$
group multiplications.
The tests in the displayed
definitions are equality
comparisons of these circuit
outputs. Combining finitely
many polynomial-size circuits
gives the asserted finite
relation circuit.
\end{proof}

\begin{theorem}[Perfect algebraic completeness of the full protocol]\label{thm:full:completeness}
For every $f$, $z$,
$v=\widetilde f(z)$, uniform
mask $R$ and honest prover
random tape, every relation
statement submitted by the
honest prover of
Definition~\ref{def:full:protocol}
has a valid hidden witness.
Consequently the complete
protocol accepts with
probability one whenever
$\Pi_{\rm rel}$ has perfect
completeness and the honest
verifier follows the stated
authenticated-query schedule.
\end{theorem}
\begin{proof}
The initial relation is
satisfied by construction of
$Q_0$ and the Merkle tree.
Lemma~\ref{lem:full:true-transition}
and Theorem~\ref{thm:full:masked-sumcheck}
establish every Sumcheck
pre- and post-challenge
relation, including the
initial claim. Lemma~\ref{lem:full:fold-invariants}
proves the fold, degree and
oracle-word equalities at
every level.
Lemma~\ref{lem:full:fold-query-complete}
gives valid authenticated
hidden witnesses for every
queried fold triple.
Lemma~\ref{lem:full:final-complete}
gives the terminal witness.
Every Pedersen commitment
opens using its actual honest
blinding and every Merkle path
reconstructs the announced root.
Thus all circuit relations
are true; perfect completeness
of the relation engine
and deterministic correctness
of the hash checks imply
acceptance with probability one.
\end{proof}

\begin{remark}[Two verification modes]
The full-relation mode checks
all committed word symbols and
their degrees inside a general
NP circuit; it has large,
non-KBFold-sized proofs but is
a complete finite reference
protocol. The sparse mode
retains the original KBFold
local-query architecture;
its proof of degree/proximity
soundness with the enlarged
budgets $D_j$ is a separate
Chapter~14 obligation.
\end{remark}
